\documentclass[11pt,a4paper]{article}
\usepackage[T1]{fontenc}
\usepackage[utf8]{inputenc}
\usepackage{lmodern}
\usepackage{microtype}
\usepackage[a4paper,margin=2.25cm]{geometry}
\usepackage{longtable,booktabs,array,calc}
\usepackage{fancyvrb}
\usepackage{xcolor}
\usepackage{hyperref}
\usepackage{xurl}
\usepackage{fancyhdr}
\usepackage{enumitem}
\usepackage{parskip}
\usepackage{titlesec}
\definecolor{ddtadark}{RGB}{30,38,48}
\definecolor{ddtablue}{RGB}{38,82,126}
\hypersetup{colorlinks=true,linkcolor=ddtablue,urlcolor=ddtablue,citecolor=ddtablue}
\pagestyle{fancy}
\fancyhf{}
\lhead{Documentation-Driven Threat Analysis}
\rhead{Research work plan R13}
\cfoot{\thepage}
\setlength{\headheight}{14pt}
\setlist[itemize]{leftmargin=1.5em,itemsep=0.2em,topsep=0.25em}
\setlist[enumerate]{leftmargin=1.7em,itemsep=0.2em,topsep=0.25em}
\titleformat{\section}{\Large\bfseries\color{ddtadark}}{}{0em}{}
\titleformat{\subsection}{\large\bfseries\color{ddtadark}}{}{0em}{}
\titleformat{\subsubsection}{\normalsize\bfseries\color{ddtadark}}{}{0em}{}
\providecommand{\tightlist}{\setlength{\itemsep}{0pt}\setlength{\parskip}{0pt}}
\setlength{\emergencystretch}{4em}
\hbadness=10000
\hfuzz=20pt
\sloppy
\begin{document}
\section{DDTA research work plan after documentation-layer closure}\label{ddta-research-work-plan-after-documentation-layer-closure}

\textbf{WORK PLAN - REVISION 13}

\textbf{Status:} active BA3 execution plan after BA3-T1 provenance/derivation lower-bound pressure testing.

\textbf{Supersedes:} Revision 12 only for forward execution state; R1-R12 remain historical research records.

\subsection{1. Fixed prior state}\label{1-fixed-prior-state}

\begin{itemize}
\tightlist
\item
  Chapters 2-4: CLOSED / FINAL for current thesis scope.
\item
  Documentation layer: CLOSED.
\item
  BA0-R systems-modeling prior-art research: CLOSED.
\item
  BA0 responsibility/non-goals: CLOSED.
\item
  BA1 minimal BAE identity ontology: CLOSED.
\item
  BA2 relation/action vocabulary: \textbf{CLOSED BY BA2-T4}.
\item
  \texttt{BAReferent} and \texttt{BAProposition}: ACCEPTED first-class semantic identity families.
\item
  ThreatForge is a case study/tooling instrument, not DDTA semantic authority.
\end{itemize}

\subsection{2. Closed BA2 dependency}\label{2-closed-ba2-dependency}

BA3 must annotate, trace and govern provenance for the already-closed BA2 proposition semantics. It must not reopen operator/role semantics without a material counterexample.

\subsection{3. BA3 objective}\label{3-ba3-objective}

BA3 defines the smallest derivation/provenance/authority mechanics that connect governed documentation to Base Analysis while preserving:

\begin{verbatim}
governed documentation = project authority
Base Analysis          = accepted analytical representation
analysis result         = method-owned / reviewable
corrective candidate    = proposed documentation change
\end{verbatim}

\subsection{4. BA3-T1 - source-to-Base-Analysis derivation and provenance lower bound}\label{4-ba3-t1---source-to-base-analysis-derivation-and-provenance-lower-bound}

\textbf{Status: COMPLETED / PROVISIONAL PASS WITH LOWER-BOUND CANDIDATE.}

BA3-T1 pressure-tested provenance target, multiplicity, source locators, baseline context, origin states, derivation basis and authority leakage across the current corpora.

\subsubsection{Result}\label{result}

\begin{verbatim}
Provenance on BAReferent                         REQUIRED
Provenance on BAProposition                      REQUIRED
Many-to-many source <-> BA lineage               REQUIRED
Immutable governed-baseline context              REQUIRED
Logical document identity + exact locator         REQUIRED
Source prose copy as BA authority                REJECTED
GROUNDED / DERIVED / DIAGNOSTIC_UNRESOLVED        REQUIRED
Derived basis separate from ultimate authority    REQUIRED
Inspectable derivation rule/rationale reference  REQUIRED CANDIDATE
Exact derivation-rule registry                    OPEN
Acceptance/review lifecycle                       NOT CLOSED BY T1
New provenance BAE identity family                NOT FORCED
Method finding / ThreatForge state as source       REJECTED
BA1 / BA2 reopen                                  NOT TRIGGERED
\end{verbatim}

\subsubsection{Active lower-bound candidate}\label{active-lower-bound-candidate}

\begin{verbatim}
BAOriginAttachment
|- targetElement            1     BAReferent | BAProposition
|- governedBaselineKey      1
|- originState              1     GROUNDED | DERIVED | DIAGNOSTIC_UNRESOLVED
|- sourceLink               0..*  -> GovernedSourceRef
|- derivationBasis          0..*  -> GovernedSourceRef | BAElementRef
`- derivationRuleRef        0..1  [required for DERIVED; exact form OPEN]

GovernedSourceRef
|- documentIdentity         1
`- locator                  1
\end{verbatim}

The contract is representation-independent metadata, not a third BAE family.

\subsection{5. Why BA3 remains open}\label{5-why-ba3-remains-open}

T1 closes no lifecycle/equivalence contract. Material unresolved questions remain:

\begin{itemize}
\tightlist
\item
  whether a BA identity is retained or replaced when governed sources change;
\item
  when an existing element becomes stale, superseded or retired;
\item
  how accepted/rejected/review state relates to origin state;
\item
  how a diagnostic is resolved/replaced after source correction;
\item
  exact derivation-rule registry/material form;
\item
  downstream change-impact and feedback lineage.
\end{itemize}

These cannot be answered by adding more source fields; they require mutation pressure.

\subsection{6. BA3-T2 - cross-baseline identity, staleness and lifecycle pressure test}\label{6-ba3-t2---cross-baseline-identity-staleness-and-lifecycle-pressure-test}

\textbf{NEXT / NOT EXECUTED BY THIS PLAN REVISION.}

BA3-T2 must use concrete governed mutations rather than inventing lifecycle vocabulary in the abstract.

It must pressure-test at least:

\begin{enumerate}
\def\labelenumi{\arabic{enumi}.}
\tightlist
\item
  facial M1 Ethernet -\textgreater{} Wi-Fi: retained abstract connectivity meaning versus changed realization proposition/source;
\item
  facial M2 externally consumed transport -\textgreater{} project-owned transport: responsibility-boundary change and affected proposition identity;
\item
  facial M3 recognition moved onto camera: retirement/supersession of transfer obligations while unrelated referents may survive;
\item
  structurally comparable order-fulfillment responsibility-boundary changes where a provider/capability moves internal/external;
\item
  whether \texttt{BAReferent} and \texttt{BAProposition} use the same retain/revise/replace rules or need distinct lifecycle rules;
\item
  when an element becomes \texttt{STALE}, \texttt{SUPERSEDED}, \texttt{RETIRED} or remains current after source change;
\item
  whether per-element \texttt{ACCEPTED} / \texttt{REJECTED} / review state is required separately from origin state;
\item
  how \texttt{DIAGNOSTIC\_UNRESOLVED} material is resolved or superseded after governed correction;
\item
  whether any mutation forces BA1/BA2 reopen rather than only BA3 lifecycle handling.
\end{enumerate}

\subsubsection{BA3-T2 exit condition}\label{ba3-t2-exit-condition}

Produce a falsifiable identity/lifecycle candidate and the smallest unresolved set. Do not close BA3 merely because the mutation examples can be manually explained.

\subsection{7. BA4 - projections}\label{7-ba4---projections}

Only after BA3 closes, define derived human and method projections from the same accepted BA semantics. No view becomes project authority.

\subsection{8. BA5 - lexical vocabulary and optional assistance}\label{8-ba5---lexical-vocabulary-and-optional-assistance}

Only after BA3/BA4 boundaries are stable, define display labels, authoring synonyms and optional assistance. NLP/LLM proposals remain candidate-producing support.

\subsection{9. BA6 - complete Base Analysis regression and closure}\label{9-ba6---complete-base-analysis-regression-and-closure}

Regress the complete BA0-BA5 design against the closed corpora and at least one structurally different holdout. BA6 may reopen only the smallest earlier responsibility on a material counterexample.

\subsection{10. Analysis envelope remains downstream}\label{10-analysis-envelope-remains-downstream}

Only after Base Analysis closure:

\begin{itemize}
\tightlist
\item
  A1 - AnalysisRecord / execution envelope;
\item
  A2 - common reviewed Finding boundary;
\item
  A3 - change/provenance integration;
\item
  formal methodology overlays and STRIDE/STRIDE-AI evaluations.
\end{itemize}

The earlier STRIDE transfer probe remains only a bounded consumer probe.

\subsection{11. Next authorized microstep}\label{11-next-authorized-microstep}

Only after the BA3-T1 package is reviewed, committed, pushed and remotely verified, execute:

\begin{quote}
\textbf{BA3-T2 - cross-baseline identity, staleness and lifecycle pressure test.}
\end{quote}

Do not start BA4, formal threat-method overlays, Common Finding schema or implementation work.

\end{document}
